\documentclass[11pt]{article}
\usepackage{amsmath,amssymb,amsthm,booktabs,longtable,array}
\usepackage[margin=1in]{geometry}

\newtheorem{lemma}{Lemma}
\newtheorem{proposition}{Proposition}
\newtheorem{definition}{Definition}
\newtheorem{remark}{Remark}
\newtheorem{corollary}{Corollary}
\newcommand{\States}{\mathcal S}
\newcommand{\Coords}{\mathcal I}
\newcommand{\Pair}{\Pi}
\newcommand{\Exit}{\mathsf{Exit}}
\newcommand{\Block}{\mathsf{Block}}

\title{Node [131] repair laboratory:\\
sequential Boolean extension and first-conflict routing}
\author{Structural-exhaustion repair sketch}
\date{15 July 2026}

\begin{document}
\maketitle

\paragraph{Status.}
This file is deliberately separate from
\texttt{proofs/erdos\_64\_eg/erdos\_64\_proof.tex}.  It freezes the exact
node-[131] defect and records successive repair candidates.  The sequential
Boolean-extension candidate below is rejected: its negative output has no
graph-local consumer.  The second iteration replaces that candidate by an
exact all-pair anchor ledger.  The third iteration identifies the missing
legal consumer: the first selected port is already a primitive capacity
token, so the blocker-free anchor fibre enters the existing non-near-cubic
primitive-carrier audit.  The source-manuscript baseline is
\begin{center}
\texttt{4e7db39e6bc8defa1a5150b912eba4fe0208d09289db4cfdddf8f8b9d2fa7680}.
\end{center}
No claim in this laboratory is an additional hypothesis of the main theorem.

\section{Defect freeze}

The incoming path in Part X is
\[
[125]\longrightarrow[126]\longrightarrow[127]\longrightarrow[128]
\longrightarrow[129]\longrightarrow[130].
\]
Node [130] forms the exact pair partition supplied by the canonical local
blocker scan.  Write its pair-local blocker-free side as
$\Pair_{\mathrm{loc}}$ and its admitted blocked side as
$\Pair_{\mathrm{blk}}^0$.  Node [131] currently receives
\[
\Coords_0:=\Coords_{\mathrm{spine}}
\qquad\text{and}\qquad
\{r_\pi:\pi\in\Pair_{\mathrm{loc}}\}.
\]
The baseline family $\Coords_0$ is independently target-testable.  The CT15
output currently established for $\Pair_{\mathrm{loc}}$ is injectivity of the
declared coordinate labels under every admitted pair-identifying quotient.

The failed implication is
\[
\begin{split}
&\text{every admitted quotient is label-injective on }
  \{r_\pi:\pi\in\Pair_{\mathrm{loc}}\}\\
&\hspace{35mm}\Longrightarrow
\text{the mixed family realizes every vector in }
\{0,1\}^{\Coords_0\cup\{r_\pi:\pi\in\Pair_{\mathrm{loc}}\}}.
\end{split} \tag{D131}
\]
This implication is false.  It is not made true by the fact that each
coordinate separately takes both values.  The three-state response relation
\[
\{(0,0),(0,1),(1,0)\}\subsetneq\{0,1\}^2
\]
has distinct, context-separable coordinate labels and omits $(1,1)$.  The
corresponding finite countermodel is kernel checked in
\texttt{StructuralExhaustion.Examples.BooleanStateEntropy}.

Consequently the node-[131] residual is the literal negation of mixed
realization:
\[
\mathsf{Missing}_{131}:\quad
\exists\varepsilon:\Coords_0\cup
  \{r_\pi:\pi\in\Pair_{\mathrm{loc}}\}\to\{0,1\},\quad
\forall s\in\States,\ \operatorname{val}(s)\ne\varepsilon . \tag{R131}
\]
The residual is not an assumption.  It is the negative branch that the
repair must consume.

\section{Frozen directed branch state}

The branch tuple at [131] is
\[
\mathfrak B_{131}=(H_0,\prec,\mathcal E,\mathcal I,
\mathcal R,\mathcal V,\mathcal Q,\mathcal A),
\]
with the following fields.
\begin{longtable}{>{\raggedright\arraybackslash}p{0.17\textwidth}
                  >{\raggedright\arraybackslash}p{0.75\textwidth}}
\toprule
Field & Value and producer \\
\midrule
$H_0$ & The selected finite simple target-avoiding graph, minimum degree at
least three, and lexicographic minimality inherited before [125]. \\
$\prec$ & The declared vertex, surplus-port, pair, blocker, and certificate
orders fixed before [130]. \\
$\mathcal E$ & All sparse exits in the manuscript definition of sparse surplus
exits are absent on this survivor. \\
$\mathcal I$ & The sparse upper envelope; the exact slack--surplus identity;
the full active excess-port family of cardinality $\sigma(G)$; all retained
activation supports and paths; the baseline family $\Coords_0$ and its
recorded deficit; the exact pair-local blocker scan and pair partition. \\
$\mathcal R$ & The unprocessed pair-response schedule together with the mixed
realization residual (R131) when it occurs. \\
$\mathcal V$ & Declared graph vertices, edge incidences, surplus ports,
activation supports, pair supports, response coordinates, admitted blockers,
capacity tokens, and proof-carrying local certificates. \\
$\mathcal Q$ & The admitted blocked-pair token payloads already queued for
[134]--[136]. \\
$\mathcal A$ & The Part-X diagram rows, the invariant ledger, the payload
registry, and the TeX--Lean index through [130]. \\
\bottomrule
\end{longtable}

No near-cubic estimate is available in this tuple.  No Type B fan-mass bound
that assumes the near-cubic spine is available either.  Facts proved only on
those branches cannot be used to repair [131].

\section{The finite local repair machine}

The complete Boolean cube must not be materialized.  The correct local object
is a compact extension operation and its first failure.

\begin{definition}[Prefix realization]
Let $c_1,\ldots,c_k$ be the ordered pair coordinates and let $\Coords_0$ be
the baseline coordinate family.  Put
\[
\Coords_i=\Coords_0\cup\{c_1,\ldots,c_i\}.
\]
A prefix constructor at stage $i$ is a function, at proof level,
\[
F_i:\bigl(\Coords_i\to\{0,1\}\bigr)\longrightarrow\States
\]
such that $\operatorname{val}(F_i(\alpha))|_{\Coords_i}=\alpha$.  The function
is not evaluated on all assignments; it is a uniform construction proved
correct for an arbitrary assignment.
\end{definition}

\begin{definition}[One-coordinate extension and first conflict]
Assume a prefix constructor $F_i$ has been retained.  The next coordinate
$c_{i+1}$ is \emph{extendible} if, for every assignment $\alpha$ on
$\Coords_i$ and every bit $b$, a locally certified operation returns a state
$s'$ satisfying
\[
\operatorname{val}(s')|_{\Coords_i}=\alpha,
\qquad
\operatorname{val}(s')(c_{i+1})=b. \tag{E}
\]
The negative branch is the lexicographically first tuple $(i+1,\alpha,b)$
for which the local extension audit cannot construct (E).  Its certificate
must be a bounded graph object checked on the declared support of
$c_{i+1}$ and the retained prefix interface; absence of a state after scanning
all labelled graphs is not an admissible certificate.
\end{definition}

\begin{lemma}[Sequential realization]
If $F_0$ realizes the baseline family and every ordered pair coordinate is
extendible in the sense of (E), then the full mixed coordinate family has a
complete Boolean-state realization.
\end{lemma}

\begin{proof}
Induct on $i$.  The base constructor is $F_0$.  Given $F_i$ and an assignment
$\beta$ on $\Coords_{i+1}$, apply the certified extension operation to
$\alpha=\beta|_{\Coords_i}$ and $b=\beta(c_{i+1})$.  The returned state has
the required values on $\Coords_i$ and on $c_{i+1}$, hence realizes $\beta$.
This defines $F_{i+1}$.  After $k$ steps, $F_k$ realizes every mixed Boolean
assignment.  The proof performs $k$ symbolic extension steps for one arbitrary
assignment; it never enumerates the $2^k$ assignments.
\end{proof}

\begin{lemma}[First-conflict dichotomy]
For the finite ordered pair schedule, either the sequential construction
reaches $F_k$, or there is a least failed extension index.  On the failure
branch the unprocessed-coordinate measure
\[
\mu=k-i
\]
strictly decreases after the failed coordinate is consumed by a registered
route.
\end{lemma}

\begin{proof}
The schedule is finite and ordered.  Scan it once.  A successful extension
increments $i$.  If an extension fails, retain the least failing index and its
local certificate.  A legal route must either close the graph or remove that
coordinate from the unprocessed schedule; in the latter case $\mu$ decreases
by one.  Hence re-entry terminates after at most $k$ successful extensions or
routed failures.
\end{proof}

This is the CT10--CT15--CT6 control flow required by the framework: CT10 names
the next coordinate and its finite response label, CT15 tests target-relative
extension, and CT6 retains the first failed local extension.  The successful
terminal supplies the exact realization contract consumed by the entropy
bound.  A failed terminal is not called a rank quotient; it is the positive
first-conflict residual.

\section{Required graph classifier}

For the repair to close, the following proposition must be proved from the
pre-[131] graph state.  It is stated here as the exact remaining proof target,
not as an assumption or a result.

\begin{quote}
\textbf{Sparse pair product-or-route target.}
Let $\Coords_i$ be a realized mixed prefix and let
$c_{i+1}=r_\pi$ be the next pair-local blocker-free response coordinate.
For an arbitrary prefix assignment $\alpha$ and requested bit $b$, exactly
one of the following certified outcomes is constructed from the declared
support data.
\begin{enumerate}
\item A state extension satisfying (E).
\item A payload satisfying the exact trigger of one of the already registered
sparse exits: direct target, target-defective quotient, proper
target-complete compression, proper or closed delocalization, or suppression
arithmetic contradiction.
\item An admitted structural blocker for $\pi$, together with its canonical
capacity token and role.  The coordinate is moved from the current free
schedule to the blocked ledger, after which the scan re-enters with the same
realized prefix and one fewer unprocessed coordinate.
\end{enumerate}
The alternatives are exhaustive; their certificate check is polynomial in
the declared supports and does not enumerate contexts, graphs, or Boolean
assignments.
\end{quote}

If this proposition is proved, the repaired node is total.  Let
$\Pair_{\mathrm{ind}}$ be the successfully extended coordinates and let
$\Pair_{\mathrm{blk}}$ be the initial admitted blockers plus all failed
coordinates routed by outcome 3.  Then
\[
\binom{\mathcal A_0}{2}
=\Pair_{\mathrm{ind}}\sqcup\Pair_{\mathrm{blk}}.
\]
The sequential lemma gives complete mixed realization for
$\Coords_0\cup\{r_\pi:\pi\in\Pair_{\mathrm{ind}}\}$, so the entropy
sandwich charges exactly $|\Pair_{\mathrm{ind}}|$.  The existing token ledger
charges exactly $|\Pair_{\mathrm{blk}}|$.  Thus the downstream pair identity
and capacity comparison are unchanged, but every coordinate reaches them
through a proved local route.

The current definitions do not prove the product-or-route target.  In
particular, pair-local absence of the listed blocker types does not define an
operation that changes $r_\pi$ while preserving an arbitrary earlier mixed
assignment.  Nor does a missing mixed vector construct a boundary mismatch,
a distinguishing outside context, a smaller representative, or one of the
current admitted blocker values.  This is precisely the missing graph lemma;
the abstract quotient-injectivity theorem cannot replace it.

\section{Audit of reuse of existing branches}

\begin{longtable}{>{\raggedright\arraybackslash}p{0.18\textwidth}
                  >{\raggedright\arraybackslash}p{0.23\textwidth}
                  >{\raggedright\arraybackslash}p{0.22\textwidth}
                  >{\raggedright\arraybackslash}p{0.27\textwidth}}
\toprule
Candidate consumer & Exact entry trigger & Dependency result & Present verdict \\
\midrule
Sparse exit: target defect & A concrete boundary-preserving raw proposal and
a compatible context with unequal target responses & Upstream of [131] and
independent & Legal only when the first-conflict classifier constructs this
payload; (R131) alone does not. \\
Sparse exit: proper compression & A target-complete proper-support quotient
with a certified strictly smaller representative & Upstream and independent &
Legal only with that representative; a missing vector does not supply one. \\
Sparse exit: delocalization & A determination certificate with its exact
minimal connected support & The proper/closed minimality argument is
independent of [131] & Legal only after constructing the determination
certificate; quotient-label injectivity supplies none. \\
Suppression exit & A retained critical suppression cycle and a chord-set
arithmetic violation & Upstream and independent & Legal only when the local
extension audit constructs that exact cycle certificate. \\
Initial blocked-pair ledger & One of the admitted blocker values and its
canonical capacity token & Nodes [134]--[144] consume it after the repaired
partition & Legal and noncircular if the new classifier proves the exact
blocker and token fields. \\
Curvature rank-drop branch & A functional target-dependence with a
determination certificate & Its intended local exits are upstream, but its
trigger is stronger than a missing Boolean vector & Illegal from (R131): the
three-state countermodel has a missing vector without a functional
coordinate dependence. \\
Type B fan-mass closure & Decorated fan or B2-overlap payload plus the
near-cubic surplus bound used by the fan-mass estimate & Depends on the spine
whose derivation passes through [131] & Circular.  A stronger local Type B
C1--C4 certificate could be used, but the generic fan-mass destination cannot. \\
Node [138] & The completed free/blocked quadratic budget or the fixed-cap
bottleneck conclusion & Descendant of [131] & Circular as a direct repair
destination. \\
HSS closure & A finite induced-$P_{13}$-free graph of minimum degree at least
three & Independent external input & Legal only if the local conflict
constructs that exact graph; (R131) does not imply it. \\
\bottomrule
\end{longtable}

Thus routing to an already closed branch is part of the proposed classifier,
but the route is legal only for sparse-exit or direct-target payloads whose
full triggers are constructed locally.  A green diagram status is not a
dependency proof.

\section{Both-sides and practicality audit}

\begin{center}
\begin{tabular}{p{0.20\textwidth}p{0.25\textwidth}p{0.25\textwidth}p{0.20\textwidth}}
\toprule
Predicate & Positive payment & Negative residual & Consumer \\
\midrule
The next response coordinate admits the certified extension (E) & Add one
coordinate to the compact prefix constructor; later charge one entropy unit &
First failed tuple $(i,\alpha,b)$ with its bounded support certificate &
Product-or-route classifier, then sparse exit or blocker ledger \\
\bottomrule
\end{tabular}
\end{center}

The schedule has at most polynomially many pair coordinates.  Each successful
step verifies one proof-carrying local operation.  Each failed step verifies
one local route and decreases the unprocessed-coordinate measure.  The
realization constructor is universally quantified Lean proof data rather than
a table of $2^k$ states.  Therefore the intended execution is polynomial in
the pair schedule and declared certificate sizes.  The only unproved item is
the graph classifier above; replacing it by a scan of all labelled graphs or
all Boolean vectors would violate the framework and is not an admissible
repair.

\section{Second repair iteration: total anchor-token accounting}

The first-conflict construction is not retained as the repair invariant.  Its
positive side is correct, but its negative side starts with an arbitrary
prefix assignment and therefore does not construct any of the finite graph
objects consumed by a sparse exit.  The repair returns to node~[130] and
tokenizes every scheduled pair before the free/blocked distinction is used for
accounting.

\begin{definition}[Canonical anchor token]
Let
\[
  p_1<\cdots<p_s,\qquad s=|\mathcal A_0|=\sigma(G),
\]
be the exact activated surplus-slot schedule.  For the canonically oriented
pair \(\pi=(p_i,p_j)\) with \(i<j\), define
\[
  \Theta_{\rm anc}(\pi):=p_i.
\]
The anchor-token universe is exactly \(\mathcal A_0\).  Its fibre at \(p_i\)
is
\[
  F_{p_i}=\{(p_i,p_j):i<j\},
\]
so every two members have the same first endpoint.  Thus each fibre is a
literal star in the scheduled-pair graph.
\end{definition}

\begin{lemma}[Exact all-pair anchor ledger]
The anchor fibres are pairwise disjoint and partition the whole scheduled
pair family.  Consequently
\[
  \binom{s}{2}=\sum_{p\in\mathcal A_0}|F_p|.
\]
The token supply is exactly \(s=\sigma(G)\), and the complete pair schedule
has quadratic size in \(s\).
\end{lemma}

\begin{proof}
Every canonically oriented pair has one first endpoint and hence belongs to
one displayed fibre.  Equality of two anchor tokens is equality of their
first endpoints, so distinct fibres are disjoint.  The cardinality identity
follows.  The token-supply identity is the already verified activated-slot
identity \(|\mathcal A_0|=\sigma(G)\).  Pair generation uses the ordered
distinct-pair iterator and is quadratic in the slot count.
\end{proof}

This lemma is kernel checked in the reusable module
\[
\texttt{StructuralExhaustion.CT9.AnchoredPairLedger}
\]
and the non-Erd\H{o}s transfer fixture
\texttt{StructuralExhaustion.Examples.CT9AnchoredPairLedger}.  The thin
problem instantiation is
\texttt{Erdos64EG.Internal.allPairAnchor\_noOvercounting}; its exact token
supply and star-fibre theorems are
\texttt{allPairAnchor\_tokenSupply\_eq\_sigma} and
\texttt{allPairAnchor\_sameAnchor}.

\begin{corollary}[Bounded anchor fibres close the numerical side]
If \(|F_p|\le L\) for every anchor token, then
\[
  \binom{\sigma(G)}2\le L\sigma(G).
\]
For \(\sigma(G)>0\) this implies \(\sigma(G)\le2L+1\), hence in
particular the near-cubic estimate needed at node~[138].
\end{corollary}

The complementary residual is now completely local and contains no Boolean
assignment:
\[
\mathsf{AnchorOverload}(p,L):\qquad |F_p|>L.
\tag{A131}
\]
CT9 constructs the exact overloaded fibre, and the reusable anchor theorem
constructs its common first endpoint.  This is strict progress over
\((R131)\): the residual is a finite star of actual activated surplus pairs.

\subsection*{Consumer audit for the overloaded star}

The current blocked-token route cannot consume (A131).  Its trigger contains
a canonical structural blocker, one of the capacity tokens attached to that
blocker, and one of the six blocker roles.  A blocker-free member of
\(F_p\) has none of those fields.  Re-labelling the anchor as a capacity token
would establish the ledger identity but would not construct the blocker or
the geometric role required by nodes [134]--[144].

The current decorated Type B route also cannot consume (A131).  Its exact
trigger contains a connected negative counted core \(Y\) in the
\(P_{13}\)-free remainder, a high decoration outside that core, a nonempty
set of distinct first neighbours, proof-carrying arms with first-entry
certificates, pairwise fan-safety, context safety, forbidden-freeness,
core-freeness, and uncompressibility.  The fact that scheduled pairs share
their first surplus slot proves none of the negative-core, first-neighbour,
or fan-safety fields.  Moreover the generic Type B bridge-mass conclusion
uses the near-cubic bound derived after node~[131], so it is not a legal
closure for this residual.

Finally, the same-token bottleneck prose currently present downstream cannot
be imported as a proof of the new route.  Its parallel and cubic-separator
cases pass from target-completeness of an identification to a proper smaller
representative.  The admissible-quotient definition requires that
representative as input; target-completeness alone does not construct it.
The all-pair ledger therefore removes the false Boolean-product implication,
but it does not by itself prove a fixed cap on blocker-free anchor fibres.

\section{Merge gate after the second iteration}

The all-pair anchor ledger is a proved reusable component, but the complete
repair still fails the leaf-totality test at (A131).  A merge would require a
graph theorem which, from the exact overloaded anchor star, constructs either
an already registered sparse-exit payload or the full decorated Type B entry
payload without using the near-cubic spine.  No current theorem has that
statement, and none of its missing payload fields follows from the anchor
ledger.  Until a different invariant closes this residual:
\begin{enumerate}
\item node [131] remains open;
\item no complete-realization hypothesis may be carried downstream;
\item the exact anchor ledger may be used as bookkeeping, but not as a
geometric closure certificate;
\item no existing closed branch may be reused without the trigger and
dependency proofs listed above.
\end{enumerate}

\section{Third repair iteration: exact route into the non-near-cubic ledger}

The overloaded anchor star is not an attempted Boolean-realization theorem.
It is already one of the three token classes of the Part-X residual.  Indeed,
the primitive carrier universe contains the selected surplus ports
\(\mathcal P_{\rm exc}\).  Thus the first endpoint of a canonically oriented
free pair is a primitive selected-port token already counted by the manuscript.

\begin{definition}[Total pair token route]
For a canonically oriented pair \(\pi=(p_i,p_j)\), \(i<j\), define
\[
 \widehat\Theta(\pi)=
 \begin{cases}
   \Theta_{\rm cap}(\pi),&\pi\in\Pi_{\rm blk},\\
   p_i,&\pi\in\Pi_{\rm free}.
 \end{cases}
\]
On the blocked side retain the canonical blocker kind at the node-[130]
dispatch and compute the existing blocker/token subtype only at nodes
[134]--[136].  On the free side use the single role
\(\mathsf{freeAnchor}\).  After retokenization the final role alphabet is
\[
 \widehat{\mathfrak R}_{\rm st}
 =\{\mathsf{blocked}(r):r\in\mathfrak R_{\rm st}\}
   \sqcup\{\mathsf{freeAnchor}\},
 \qquad |\widehat{\mathfrak R}_{\rm st}|\le4\cdot6+1=25.
\]
The free role occurs only with the primitive selected-port subtype.
\end{definition}

\begin{lemma}[Exact total-pair route]
The fibres of \((\widehat\Theta,\widehat\rho)\) form a disjoint partition of
the complete scheduled pair family.  In particular,
\[
 \binom{|\mathcal A_0|}{2}
 =\sum_{C\in\{W,R,P\}}
   \sum_{t\in\mathfrak T_C}
   \sum_{r\in\widehat{\mathfrak R}_{\rm st}}\widehat\ell(t,r).
 \tag{T131}
\]
Every member of a \((p,\mathsf{freeAnchor})\)-fibre is blocker-free, has
first endpoint exactly \(p\), and retains the shortest connector and both
activated demand supports selected at node~[130].
\end{lemma}

\begin{proof}
The canonical blocker scan gives the decidable disjoint partition
\(\Pi_{\rm blk}\sqcup\Pi_{\rm free}\).  A blocked pair is sent through the
existing single-valued map \(\Theta_{\rm cap}\) and its existing role map.  A
free pair has a unique first endpoint in the ordered pair schedule and is sent
to that endpoint with the distinct role \(\mathsf{freeAnchor}\).  Hence the
two images cannot overlap, and within either side equality of the token--role
label defines the displayed fibres.  This proves (T131).  For the last
assertion, membership in the free-anchor fibre unfolds to the two equalities
\(p_i=p\) and \(\widehat\rho(\pi)=\mathsf{freeAnchor}\).  The latter is
definitionally the negative result of the canonical blocker decision, while
the connector and its two support inclusions are the data already retained
for that free pair.
\end{proof}

This is kernel checked in three layers.  The reusable product partition and
local-cap aggregation are
\texttt{StructuralExhaustion.CT9.TokenRoleLedger}; the non-Erd\H{o}s transfer
fixture is \texttt{StructuralExhaustion.Examples.CT9TokenRoleLedger}; and the
graph-owned semantic route is
\texttt{StructuralExhaustion.Graph.SurplusPairTokenRouting}.  Its Erd\H{o}s
endpoint is
\texttt{Erdos64EG.Internal.exists\_verifiedAllPairTokenRoutingPrefix}.
The executable node-[130] dispatch uses five roles---the four admitted
structural blocker kinds and \(\mathsf{freeAnchor}\)---because profile and
target mismatches leave as raw audit exits and a blocked token subtype is not
invented before the downstream capacity-token map.  The final combined ledger
uses the 25-role alphabet above.

\subsection*{Exact consumer and dependency audit}

For a free-anchor fibre the token \(p\) lies in
\(\mathcal P_{\rm exc}\subseteq\mathfrak T_{\rm prim}\).  Its fibre is a
literal role-homogeneous star in the pair graph, centered at \(p\).  Therefore
its exact consumer is node~[143], the primitive-carrier geometric audit.  The
support passed to that audit is the retained pair connector support, not a
newly searched subgraph.  In the routing-germ definition the token root is
the buffer \(x(p)\), exactly the already specified root of a primitive
selected-port token.

This handoff is noncircular.  Nodes [125]--[130] produce the graph, activated
slots, blocker decision, and free-pair connectors.  The route uses only those
objects and the fixed primitive-token inclusion.  Node [143] does not become
an input hypothesis: it is the next residual audit.  In particular, no
near-cubic estimate or Type-B fan-mass estimate is transported into the
route.  The later alternatives of the geometric audit may produce a sparse
exit, Type-B handoff data, or fixed caps; only the fixed-cap outcome is then
allowed to produce node~[138].

The blocked branch remains unchanged: its canonical first-hit record, with
the absence of all earlier candidates, enters node~[134], where the existing
capacity-token map computes the token subtype.  Thus the total route neither
adds a blocker hypothesis to a free pair nor assigns a fictitious subtype to
a blocked pair.

\section{Merge gate after the third iteration}

The node-[131] repair passes the local leaf-totality test.  Every pair has one
of two exact outputs: a canonical blocker handoff to node~[134], or a
free-anchor primitive-token handoff to node~[143].  The false implication
(D131) and the mixed-realization residual (R131) are no longer used by the
Part-X accounting route.  The frontier after this repair is the downstream
capacity-token and primitive geometric audits themselves; their verification
status is independent of the completed node-[131] dispatch.

\end{document}
